% Generated draft: requires mathematical and source review.
\documentclass[11pt]{article}
\usepackage[margin=25mm]{geometry}
\usepackage{fontspec}
\setmainfont{Latin Modern Roman}
\usepackage{amsmath,amssymb,mathtools,unicode-math}
\setmathfont{Latin Modern Math}
\usepackage{enumitem,needspace,xurl,hyperref}
\setlength{\parindent}{0pt}
\setlength{\parskip}{5pt}
\setlength{\emergencystretch}{3em}
\newcommand{\N}{\mathbb{N}}
\newcommand{\Z}{\mathbb{Z}}
\newcommand{\Q}{\mathbb{Q}}
\newcommand{\R}{\mathbb{R}}
\newcommand{\C}{\mathbb{C}}
\newcommand{\F}{\mathbb{F}}
\newcommand{\A}{\mathbb{A}}
\newcommand{\PP}{\mathbb{P}}
\newcommand{\E}{\mathbb{E}}
\newcommand{\eps}{\varepsilon}
\newcommand{\dd}{\,\mathrm{d}}
\begin{document}
\section*{UnsolvedMath ID: 800147}
\textbf{Source problem code:} catalogue entry
\subsection{Definitions and mathematical statement}
The supplied source asks for numerical computations of further prime number races with at least three competitors. The following makes explicit a conventional interpretation; the source excerpt does not specify the density convention or the hypotheses under which the relevant densities exist.

Let \(q\geq 2\) be an integer, and let \(r\geq 3\). Choose integers \(a_1,\ldots,a_r\) representing distinct residue classes modulo \(q\), with \(\gcd(a_i,q)=1\) for every \(i\). Thus \(r\leq\varphi(q)\), where \(\varphi(q)\) is the number of invertible residue classes modulo \(q\). For real \(x\geq 2\), define
\[
\pi(x;q,a)=\#\{p\leq x:p\text{ is prime and }p\equiv a\pmod q\}.
\]
For the ordered list \((a_1,\ldots,a_r)\), define the strict-order race set
\[
\mathcal R(q;a_1,\ldots,a_r)
=\{x\geq 2:\pi(x;q,a_1)>\pi(x;q,a_2)>\cdots>\pi(x;q,a_r)\}.
\]
Its logarithmic density, when the limit exists, is
\[
\delta(q;a_1,\ldots,a_r)
=\lim_{X\to\infty}\frac{1}{\log X}
\int_2^X \mathbf 1_{\mathcal R(q;a_1,\ldots,a_r)}(t)\,\frac{\dd t}{t},
\]
where \(\mathbf 1_S\) denotes the indicator of a set \(S\).

\textbf{Computational problem, under this interpretation.} Choose further ordered lists of at least three distinct reduced residue classes and numerically estimate their race densities \(\delta(q;a_1,\ldots,a_r)\). Identify the modulus and ordered competitors for each computation, and state any hypotheses used to justify existence of the density or applicability of a numerical formula. The supplied excerpt imposes no particular modulus, list of competitors, number of examples, or numerical accuracy requirement.

The source reports that Feuerverger and Martin [FeM] computed densities for some three-way races and for four-way races modulo \(5\), \(8\), and \(12\), and supplied a general density formula for any number of competitors. It asks for further computations and describes subsequent computations with three or more competitors as few, if any. The excerpt does not enumerate the earlier three-way examples or provide the formula and its hypotheses, so it does not by itself determine which proposed examples would be new. It also reports that Fiorilli and Martin [FiM] determined the \(117\) most biased two-way races. These are supplied background claims, not independently verified conclusions. The reduced-residue-class model, strict ordering, and logarithmic-density convention above require confirmation against the full source.
\subsection{Short English statement}
Compute additional examples of the densities of prime number races with three or more competitors. A conventional formulation compares prime counts in distinct reduced residue classes modulo a common integer and estimates the logarithmic density of a specified strict ordering. The supplied excerpt leaves the precise density convention, existence hypotheses, examples, and numerical precision unspecified; this statement-only draft requires human mathematical review.
\subsection{Sources}
\begin{itemize}
\item Local UnsolvedMath registry: permanent ID \(800147\), exact source problem code \texttt{catalogue entry}, title \emph{Comparative Prime Number Theory --- Problem 1}. The supplied record labels the problem \texttt{open}; that label is source metadata, not an independently verified mathematical assessment. The record was created and updated on 2026-10-16.
\item Hamieh, Kadiri, Martin, and Ng, \emph{Comparative Prime Number Theory Problem List} (2024), source item \emph{Problem 1 (\#16)}. Supplied URL: \url{https://arxiv.org/abs/2407.03530}. The record reports access on 2026-10-16 and extraction by \texttt{arxiv\_latex}; these are supplied provenance statements, not a claim of access or verification for this draft. The item is attributed within the excerpt to Nathan Ng, University of Lethbridge; the registry lists Hamieh, Kadiri, Martin, and Ng under \texttt{proposed\_by}, with proposed year 2024.
\item Feuerverger and Martin [FeM], cited in the supplied excerpt for earlier numerical race densities and a general formula. No full bibliographic entry, direct link, formula, or hypotheses for this reference are supplied.
\item Fiorilli and Martin [FiM], cited in the supplied excerpt for determining the \(117\) most biased two-way races. No full bibliographic entry, direct link, or definition of the ranking criterion is supplied.
\end{itemize}
\end{document}
